\documentclass[a4paper,11pt]{article}
\usepackage[utf8]{inputenc}
\usepackage[T1]{fontenc}
\usepackage[french]{babel}
\usepackage{geometry}
\geometry{left=2cm,right=2cm,top=1.5cm,bottom=2cm}
\usepackage{amsmath,amssymb}
\usepackage{enumitem}
\usepackage{array}
\usepackage{booktabs}
\usepackage{xcolor}
\usepackage{tcolorbox}
\usepackage{fancyhdr}
\usepackage{tikz}
\usepackage{lastpage}

% Couleurs
\definecolor{bleupro}{HTML}{0969da}
\definecolor{orange}{HTML}{d97706}
\definecolor{vertpro}{HTML}{2f855a}
\definecolor{rouge}{HTML}{c53030}
\definecolor{gris}{HTML}{718096}

% En-tête
\pagestyle{fancy}
\fancyhf{}
\lhead{\small CCF Mathématiques — Terminale Bac Pro}
\rhead{\small Probabilités \& Polynômes degré 3}
\cfoot{\thepage/\pageref{LastPage}}
\renewcommand{\headrulewidth}{0.5pt}

% Boîtes
\tcbuselibrary{skins,breakable}

\newtcolorbox{formulairebox}{
  colback=yellow!5,
  colframe=orange,
  title={\textbf{FORMULAIRE — Rappels de cours}},
  fonttitle=\large\bfseries,
  breakable
}

\newtcolorbox{contextebox}{
  colback=blue!3,
  colframe=bleupro,
  leftrule=4pt,
  rightrule=0pt,
  toprule=0pt,
  bottomrule=0pt,
  boxsep=4pt,
  left=8pt
}

\newtcolorbox{reponsebox}[1][4cm]{
  colback=white,
  colframe=gris!40,
  boxrule=0.5pt,
  arc=3pt,
  height=#1,
  valign=top,
  before skip=4pt,
  after skip=8pt
}

\newcommand{\pts}[1]{\hfill\textcolor{gris}{\small\textbf{#1}}}
\newcommand{\comp}[1]{\colorbox{blue!10}{\scriptsize\textbf{#1}}}

\setlength{\parindent}{0pt}
\setlength{\parskip}{6pt}

\begin{document}

% ===== PAGE DE GARDE =====
\begin{center}
{\LARGE\bfseries\color{bleupro} Contrôle en Cours de Formation}\\[8pt]
{\large Mathématiques — Terminale Bac Pro — ERA-TMA / ICCER}\\[12pt]
{\Large\bfseries Probabilités \& Fonctions polynômes de degré 3}\\[20pt]

\begin{tabular}{|p{6cm}|p{6cm}|}
\hline
\textbf{Durée :} 45 minutes & \textbf{Barème :} /10 points \\
\hline
\textbf{Documents :} Calculatrice NumWorks autorisée & \textbf{Formulaire :} ci-dessous \\
\hline
\textbf{Nom :} \dotfill & \textbf{Prénom :} \dotfill \\
\hline
\textbf{Classe :} \dotfill & \textbf{Date :} \dotfill \\
\hline
\end{tabular}
\end{center}

\vspace{10pt}

% ===== FORMULAIRE =====
\begin{formulairebox}

\subsection*{Probabilités}

\begin{tabular}{>{\bfseries}p{5cm} p{8.5cm}}
\toprule
\textbf{Notion} & \textbf{Formule / Règle} \\
\midrule
Probabilité d'un chemin & Multiplier les probabilités le long des branches \\[4pt]
Probabilités totales & $P(B) = P(A) \times P_A(B) + P(\overline{A}) \times P_{\overline{A}}(B)$ \\[4pt]
Probabilité conditionnelle & $P_A(B) = \dfrac{P(A \cap B)}{P(A)}$ \quad avec $P(A) \neq 0$ \\[6pt]
Événement contraire & $P(\overline{A}) = 1 - P(A)$ \\[4pt]
Indépendance & $A$ et $B$ indépendants $\Leftrightarrow$ $P(A \cap B) = P(A) \times P(B)$ \\[4pt]
\bottomrule
\end{tabular}

\smallskip
\textbf{Arbre pondéré :} la somme des probabilités des branches issues d'un même nœud vaut toujours 1.

\subsection*{Fonctions polynômes de degré 3}

\begin{tabular}{>{\bfseries}p{5cm} p{8.5cm}}
\toprule
\textbf{Notion} & \textbf{Formule / Règle} \\
\midrule
Forme générale & $f(x) = ax^3 + bx^2 + cx + d$ \quad ($a \neq 0$) \\[4pt]
Dérivée & $f'(x) = 3ax^2 + 2bx + c$ \\[4pt]
Sens de variation & $f'(x) > 0 \Rightarrow f$ croissante \quad $f'(x) < 0 \Rightarrow f$ décroissante \\[4pt]
Extremum local & Si $f'$ change de signe en $x_0$ : extremum en $(x_0\,;\,f(x_0))$ \\[4pt]
Nombre de solutions & Exploiter le tableau de variations pour $f(x) = c$ \\[4pt]
\bottomrule
\end{tabular}

\smallskip
\textbf{Attention :} $f'(x_0) = 0$ ne suffit pas pour conclure à un extremum (ex : fonction cube en $x = 0$).

\end{formulairebox}

\newpage

% ================================================================
\section*{Exercice 1 — Contrôle qualité en menuiserie \pts{5 points}}
% ================================================================

\begin{contextebox}
Un atelier de menuiserie fabrique des panneaux d'agencement. Chaque panneau passe deux contrôles successifs.

\textbf{Contrôle 1 — Dimensions :} 8\% des panneaux ont un défaut de dimensions.

\textbf{Contrôle 2 — Surface :}
\begin{itemize}[nosep]
  \item Si le panneau a un défaut de dimensions, la probabilité d'avoir aussi un défaut de surface est 0,3.
  \item Si le panneau n'a pas de défaut de dimensions, la probabilité d'avoir un défaut de surface est 0,05.
\end{itemize}

On note : $D$ = ``le panneau a un défaut de dimensions'' et $S$ = ``le panneau a un défaut de surface''.
\end{contextebox}

\vspace{6pt}

\textbf{Q1} \comp{APP} — Traduire les données de l'énoncé en probabilités : donner les valeurs de $P(D)$, $P(\overline{D})$, $P_D(S)$ et $P_{\overline{D}}(S)$. \pts{1 pt}

\begin{reponsebox}[3cm]
\dotfill

\dotfill

\dotfill
\end{reponsebox}

\textbf{Q2} \comp{REA} — Compléter l'arbre de probabilités ci-dessous. \pts{1 pt}

\begin{center}
\begin{tikzpicture}[scale=1, every node/.style={font=\small}]
  % Racine
  \node[circle, fill=black, inner sep=1.5pt] (R) at (0,0) {};
  % Niveau 1
  \node (D) at (3,1.5) {$D$};
  \node (Db) at (3,-1.5) {$\overline{D}$};
  \draw[-] (R) -- (D) node[midway, above, sloped] {\dotfill};
  \draw[-] (R) -- (Db) node[midway, below, sloped] {\dotfill};
  % Niveau 2 depuis D
  \node (DS) at (6.5,2.3) {$S$};
  \node (DSb) at (6.5,0.7) {$\overline{S}$};
  \draw[-] (D) -- (DS) node[midway, above, sloped] {\dotfill};
  \draw[-] (D) -- (DSb) node[midway, below, sloped] {\dotfill};
  % Niveau 2 depuis Db
  \node (DbS) at (6.5,-0.7) {$S$};
  \node (DbSb) at (6.5,-2.3) {$\overline{S}$};
  \draw[-] (Db) -- (DbS) node[midway, above, sloped] {\dotfill};
  \draw[-] (Db) -- (DbSb) node[midway, below, sloped] {\dotfill};
\end{tikzpicture}
\end{center}

\textbf{Q3} \comp{REA} — Calculer $P(D \cap S)$ et $P(\overline{D} \cap S)$. \pts{1 pt}

\begin{reponsebox}[3cm]
\dotfill

\dotfill

\dotfill
\end{reponsebox}

\textbf{Q4} \comp{ANA} — En utilisant la formule des probabilités totales, calculer la probabilité qu'un panneau pris au hasard ait un défaut de surface. \pts{1 pt}

\begin{reponsebox}[3cm]
\dotfill

\dotfill

\dotfill
\end{reponsebox}

\textbf{Q5} \comp{ANA} — Sur une production de 500 panneaux, combien peut-on s'attendre à avoir avec un défaut de surface ? Le responsable qualité considère que le taux est acceptable si moins de 5\% des panneaux sont défectueux en surface. Est-ce le cas ? \pts{1 pt}

\begin{reponsebox}[3cm]
\dotfill

\dotfill

\dotfill
\end{reponsebox}

\newpage

% ================================================================
\section*{Exercice 2 — Optimisation du volume d'un bac de rangement \pts{5 points}}
% ================================================================

\begin{contextebox}
Un menuisier fabrique un bac de rangement ouvert à partir d'une plaque carrée de tôle de 30 cm de côté. Il découpe un carré de côté $x$ (en cm) à chaque coin, puis plie les bords vers le haut.

Le volume du bac obtenu est :
\[V(x) = x(30 - 2x)^2 = 4x^3 - 120x^2 + 900x \quad \text{pour } x \in \;]0\,;\,15[\]
\end{contextebox}

% Schéma plaque
\begin{center}
\begin{tikzpicture}[scale=0.06]
  % Plaque
  \draw[fill=blue!5, thick] (0,0) rectangle (30,30);
  % Carrés découpés
  \draw[fill=red!10, thick, dashed] (0,0) rectangle (5,5);
  \draw[fill=red!10, thick, dashed] (25,0) rectangle (30,5);
  \draw[fill=red!10, thick, dashed] (0,25) rectangle (5,30);
  \draw[fill=red!10, thick, dashed] (25,25) rectangle (30,30);
  % Lignes de pliage
  \draw[dashed, blue!60] (5,0) -- (5,30);
  \draw[dashed, blue!60] (25,0) -- (25,30);
  \draw[dashed, blue!60] (0,5) -- (30,5);
  \draw[dashed, blue!60] (0,25) -- (30,25);
  % Cotes
  \node[below] at (15,0) {\small 30 cm};
  \node[left] at (0,15) {\small 30 cm};
  \node at (2.5,2.5) {\small $x$};
  \node at (15,15) {\small $30-2x$};
\end{tikzpicture}
\end{center}

\textbf{Q1} \comp{REA} — \textbf{À la calculatrice (NumWorks) ou sur GeoGebra}, tracer la courbe de $V$ sur l'intervalle $[0\,;\,15]$. Reproduire ci-dessous l'allure de la courbe obtenue. Placer le maximum. \pts{1 pt}

\begin{center}
\begin{tikzpicture}[scale=0.48]
  % Quadrillage
  \draw[gray!25, very thin] (0,0) grid[step=1] (16,14);
  \draw[gray!50, thin] (0,0) grid[step=5] (16,14);
  % Axes
  \draw[->, thick] (0,0) -- (16.5,0) node[right] {$x$ (cm)};
  \draw[->, thick] (0,0) -- (0,14.5) node[above] {$V$ (cm$^3$)};
  % Graduations x
  \foreach \x in {1,2,...,15} {
    \draw (\x,0) -- (\x,-0.15);
  }
  \foreach \x in {5,10,15} {
    \node[below] at (\x,-0.3) {\small\x};
  }
  % Graduations y (1 unité = 500 cm³)
  \foreach \y in {1,2,...,13} {
    \draw (0,\y) -- (-0.15,\y);
  }
  \node[left] at (-0.3,2) {\small 500};
  \node[left] at (-0.3,4) {\small 1000};
  \node[left] at (-0.3,6) {\small 1500};
  \node[left] at (-0.3,8) {\small 2000};
  \node[left] at (-0.3,10) {\small 2500};
  \node[left] at (-0.3,12) {\small 3000};
  % Origine
  \node[below left] at (0,0) {\small 0};
\end{tikzpicture}
\end{center}

\textbf{Q2} \comp{APP} — \textbf{Par lecture graphique}, donner une valeur approchée de $x$ pour laquelle le volume semble maximal. Lire la valeur approchée de ce maximum. \pts{1 pt}

\begin{reponsebox}[2.5cm]
\dotfill

\dotfill
\end{reponsebox}

\textbf{Q3} \comp{REA} — Calculer la dérivée $V'(x)$. Vérifier que $V'(x) = 12(x - 5)(x - 15)$. En déduire les solutions de $V'(x) = 0$ sur $]0\,;\,15[$. \pts{1 pt}

\begin{reponsebox}[3.5cm]
\dotfill

\dotfill

\dotfill

\dotfill
\end{reponsebox}

\textbf{Q4} \comp{ANA} — Compléter le tableau de variations de $V$ sur $]0\,;\,15[$. Calculer la valeur exacte du maximum. Ce résultat confirme-t-il la lecture graphique de Q2 ? \pts{1 pt}

\begin{center}
\renewcommand{\arraystretch}{1.8}
\begin{tabular}{|>{\centering\arraybackslash}m{2.5cm}|>{\centering\arraybackslash}m{1.5cm}|>{\centering\arraybackslash}m{2cm}|>{\centering\arraybackslash}m{1.5cm}|>{\centering\arraybackslash}m{2cm}|>{\centering\arraybackslash}m{1.5cm}|}
\hline
$x$ & $0$ & & \ldots\ldots & & $15$ \\
\hline
Signe de $V'(x)$ & & \ldots\ldots & $0$ & \ldots\ldots & \\
\hline
& & & & & \\
Variations de $V$ & \ldots\ldots & & \ldots\ldots & & \ldots\ldots \\
& & & & & \\
\hline
\end{tabular}
\renewcommand{\arraystretch}{1}
\end{center}

\vspace{4pt}
Valeur exacte du maximum : $V(\ldots\ldots) = $ \dotfill

\vspace{4pt}
Confirmation lecture graphique : \dotfill

\textbf{Q5} \comp{COM} — Conclure : donner la valeur de $x$ qui maximise le volume, le volume maximal, et les dimensions du bac (longueur, largeur, hauteur). Le menuisier a besoin d'un bac d'au moins 1\,500 cm$^3$. La contrainte est-elle respectée ? \pts{1 pt}

\begin{reponsebox}[4cm]
\dotfill

\dotfill

\dotfill

\dotfill

\dotfill
\end{reponsebox}

\vfill

\begin{center}
\begin{tabular}{|c|c|c|c|}
\hline
\textbf{Exercice} & \textbf{Thème} & \textbf{Questions} & \textbf{Points} \\
\hline
1 & Probabilités (arbre, probas totales) & Q1 à Q5 & /5 \\
\hline
2 & Polynôme degré 3 (dérivée, variations, optimisation) & Q1 à Q5 & /5 \\
\hline
\multicolumn{3}{|r|}{\textbf{Total}} & \textbf{/10} \\
\hline
\end{tabular}
\end{center}

\end{document}
